%==============================================================================
% CHAPTER 6: Orgel's Paradox and the Thermodynamic Inevitability of Life
%==============================================================================

\chapter{Orgel's Paradox and the Thermodynamic Inevitability of Life}
\label{chap:orgels_paradox}

\epigraph{%
  The origin of life is not a problem to be solved but a phenomenon to be
  understood---and the understanding begins when we ask not what came first,
  but what could not \emph{not} have come first.%
}{after Leslie Orgel, \textit{The Origins of Life} (1973)}

\begin{chapterbox}
Leslie Orgel formulated the foundational circularity of life's origin: genes
require enzymes for replication; enzymes require genes for synthesis; both
require metabolic energy; and metabolism requires enzymatic catalysis. Three
traditional escape routes---RNA world, metabolism-first, lipid world---each
privilege one component and restore the circularity at a higher level. This
chapter presents a fourth resolution, grounded in partition theory:
\emph{partitioning precedes and generates all three}. Four theorems establish
the case. The Thermodynamic Inevitability Theorem proves $\Delta G_{\mathrm{partition}} < 0$
wherever chemical potential gradients exist. The Probability Separation Theorem
shows that partition-first scenarios exceed information-first scenarios by the
categorical factor $10^{144}$. The Autocatalysis Theorem establishes that
partition operations which intensify their own gradient are self-sustaining.
The Chirality Selection Theorem derives universal homochirality from electron
spin--orbit coupling at chiral partition boundaries, without invoking any
external chiral field. The circularity of Orgel's paradox is dissolved: all
three components of life---information, catalysis, and metabolism---arise
simultaneously and independently from a single pre-biotic operation, the
partition.
\end{chapterbox}

%------------------------------------------------------------------------------
\section{The Paradox: A Genuine Impossibility}
\label{sec:orgel_paradox_statement}
%------------------------------------------------------------------------------

In the late 1960s, Leslie Orgel articulated what is perhaps the most fundamental
obstacle to understanding life's origin. Biological systems require a mutually
dependent triad:

\begin{enumerate}[label=(\Roman*)]
  \item \textbf{Genetic information}: DNA (or RNA) encodes the sequences of
        all enzymes. But DNA replication requires polymerases---enzymes.
  \item \textbf{Enzymatic catalysis}: Enzymes perform the chemical work of
        life, including the replication of DNA. But enzymes are encoded in
        DNA and translated by ribosomes, themselves assemblies of RNA and
        protein.
  \item \textbf{Metabolic energy}: Replication and translation are endergonic.
        ATP is the universal energy currency. ATP synthesis by ATP synthase
        requires a proton gradient, which the electron transport chain
        maintains---a chain of enzymes encoded in DNA.
\end{enumerate}

The circularity is exact: (I) requires (III), (III) requires (II), and (II)
requires (I). This is not a conceptual puzzle but a logical impossibility.
If everything in a living cell depends on everything else in that cell, then
the \emph{origin} of any one component requires the prior existence of all others.
No linear historical account can begin.

\begin{definition}[Orgel's Circle]
\label{def:orgels_circle}
Orgel's circle is the set of mutual logical dependencies:
\begin{equation}
  \mathrm{Gene} \xrightarrow{\text{encodes}} \mathrm{Enzyme} \xrightarrow{\text{synthesizes}}
  \mathrm{Gene}, \qquad
  \mathrm{Enzyme} \xrightarrow{\text{requires}} \mathrm{Metabolism} \xrightarrow{\text{requires}}
  \mathrm{Enzyme}.
  \label{eq:orgels_circle}
\end{equation}
No element of the circle has an independent origin within the circle's own logic.
\end{definition}

\subsection{The Three Standard Escape Attempts and Their Failures}

\subsubsection{RNA World}

Walter Gilbert's RNA world hypothesis \citep{gilbert1986} proposes that RNA
preceded both DNA and protein, because RNA can simultaneously carry sequence
information (like DNA) and catalyse reactions (like protein enzymes, through
ribozyme activity). If one molecule occupies both roles, the gene/enzyme
half of the circle is broken.

\textbf{The residual problem}: Even granting ribozyme activity, the RNA world
requires the \emph{de novo} spontaneous formation of a functional self-replicating
RNA molecule from random nucleotide chemistry. The sequence space is astronomically
large; we will quantify the probability precisely in Section~\ref{sec:probability_separation}.
More fundamentally, the RNA world still requires metabolic energy and
compartmentalisation---it dissolves one arc of the circle but restores the
full circle at the level of the first ribozyme's own energetics and boundary.

\subsubsection{Metabolism-First}

G\"{u}nter W\"{a}chtersh\"{a}user's iron--sulfur world \citep{wachtershauser1988}
proposes that autocatalytic chemical cycles on mineral surfaces preceded genetic
information. The metabolic network generates the chemical energy and molecular
building blocks that later gave rise to genetic molecules.

\textbf{The residual problem}: A spontaneous autocatalytic network requires a
specific set of coupled reactions each catalysing the next. Without sequence-based
inheritance, the network composition drifts irreversibly toward thermodynamic
equilibrium (dead chemistry). Metabolism without information is transient; it
cannot be reproduced or evolved. The metabolism-first account dissolves the
metabolic arc of Orgel's circle but cannot explain the origin of informational
specificity.

\subsubsection{Lipid World}

The lipid world proposal \citep{segre2001} holds that self-replicating
amphiphilic vesicles arose first, providing compartmentalisation and a primitive
form of reproduction without genetic information. Vesicles can grow, divide, and
encapsulate other molecules.

\textbf{The residual problem}: Vesicle self-assembly provides no catalytic
function and no mechanism for sequence-based inheritance. Division of a lipid
vesicle distributes whatever it contains but does not reproduce specific
chemical compositions. Compartmentalisation is a consequence of closed surfaces,
not a source of biological specificity.

\subsection{The Common Structural Failure}

All three escape routes share a common failure of framing: they operate at
the \emph{molecular} level and attempt to identify the ``first'' molecular
component, then show how the others follow. The partition theory diagnosis is
that this is the wrong level. The three components of life---information,
catalysis, metabolism---are not primitive. They are all manifestations of the
same underlying operation: the \emph{partition}. Partitioning is the operation
that \emph{precedes} all three, generates all three simultaneously, and
therefore dissolves the circle without privileging any component.

%------------------------------------------------------------------------------
\section{Partition Operations as Pre-Biochemical Phenomena}
\label{sec:prebiochem_partitioning}
%------------------------------------------------------------------------------

\begin{definition}[Chemical Partition Operation]
\label{def:chemical_partition}
A \emph{chemical partition operation} on a molecular system is a physical
process that divides the accessible chemical state space $\Cat$ into two or
more non-overlapping regions with the following properties:
\begin{enumerate}[label=(\alph*)]
  \item \textbf{Selectivity}: molecules in distinct regions have different
        accessible reaction pathways.
  \item \textbf{Stability}: the partition boundary $\partial\Cat$ is maintained
        against thermal fluctuations for at least one chemical timescale.
  \item \textbf{Gradient coupling}: a non-zero chemical potential gradient
        $|\nabla\mu| > 0$ is maintained across the partition boundary.
\end{enumerate}
\end{definition}

\begin{example}[Iron--Sulfur Mineral Surface as Chemical Partition]
\label{ex:iron_sulfur_partition}
Iron--sulfur minerals (mackinawite FeS, pyrite FeS$_2$, greigite Fe$_3$S$_4$),
formed spontaneously at Hadean hydrothermal vents from dissolved Fe$^{2+}$ and
H$_2$S, are semiconductors with bandgaps $E_g \approx 0.1$--$0.95\,\mathrm{eV}$.
This band structure creates a chemical partition operation:
\begin{itemize}
  \item Conduction-band electrons (above the Fermi level) drive reduction reactions.
  \item Valence-band states (below the Fermi level) receive electrons from
        oxidation reactions.
  \item The bandgap maintains selectivity between the two chemical regimes.
  \item The redox potential $E^{\circ} \approx -0.5$ to $+0.3\,\mathrm{V}$
        (vs.\ NHE) provides the chemical potential gradient.
\end{itemize}
The surface Fe--S clusters are the direct structural ancestors of the [2Fe-2S]
and [4Fe-4S] cofactors at the active sites of all ferredoxins, hydrogenases,
and nitrogenases---the oldest family of protein cofactors in biochemistry.
\end{example}

%------------------------------------------------------------------------------
\section{Theorem: Thermodynamic Inevitability of Partition Operations}
\label{sec:thermodynamic_inevitability}
%------------------------------------------------------------------------------

\begin{theorem}[Thermodynamic Inevitability]
\label{thm:thermodynamic_inevitability}
In any system with a non-zero chemical potential gradient $\Delta\mu > 0$,
partition operations have negative Gibbs free energy:
\begin{equation}
  \Delta G_{\mathrm{partition}} < 0.
  \label{eq:dg_partition_negative}
\end{equation}
Partition operations therefore occur spontaneously wherever chemical potential
gradients exist.
\end{theorem}

\begin{proof}
A partition operation creates the boundary $\partial\Part$ between regions of
different chemical potential. The free energy change has two independent
contributions.

\noindent\textbf{(1) Categorical entropy gain.} Before partitioning, the system
has microstate count $\Omega_{\mathrm{before}}$. After partitioning, the two
regions are distinguished; molecules in each region have different accessible
states. By the multiplicativity of microstates for independent subsystems,
$\Omega_{\mathrm{after}} = \Omega_1 \times \Omega_2 > \Omega_{\mathrm{before}}$
(the partition creates new categorical distinctions, increasing the total state
count). Therefore:
\begin{equation}
  \Delta S_{\mathrm{cat}} = \kB \ln\!\left(\frac{\Omega_{\mathrm{after}}}{\Omega_{\mathrm{before}}}\right) > 0.
  \label{eq:ds_cat_partition}
\end{equation}

\noindent\textbf{(2) Gradient-driven enthalpy decrease.} The gradient $\Delta\mu > 0$
drives $\Delta N$ molecules across the new partition boundary from high to low
chemical potential. The free energy extracted from this transfer is
$W_{\mathrm{gradient}} = \Delta\mu \cdot \Delta N > 0$, contributing:
\begin{equation}
  \Delta H_{\mathrm{partition}} = -\Delta\mu \cdot \Delta N < 0.
  \label{eq:dh_partition}
\end{equation}

\noindent\textbf{Combining.} The total Gibbs free energy change is:
\begin{equation}
  \Delta G_{\mathrm{partition}} = \Delta H_{\mathrm{partition}} - T\Delta S_{\mathrm{cat}}
  = -\Delta\mu \cdot \Delta N - T\kB \ln\!\left(\frac{\Omega_{\mathrm{after}}}{\Omega_{\mathrm{before}}}\right).
  \label{eq:dg_full}
\end{equation}
Both terms are negative under the stated conditions ($\Delta\mu > 0$, $\Delta N > 0$,
$\Omega_{\mathrm{after}} > \Omega_{\mathrm{before}}$). Therefore $\Delta G_{\mathrm{partition}} < 0$.
\qed
\end{proof}

\begin{corollary}[Inevitability on Early Earth]
\label{cor:early_earth_inevitability}
On early Earth, chemical potential gradients existed at: (a)~hydrothermal vents
producing H$_2$ and reduced sulfur against the oxidised ocean;
(b)~UV-irradiated surface waters; (c)~iron--sulfur mineral surfaces with
semiconductor properties in contact with gradient-bearing fluids. By
Theorem~\ref{thm:thermodynamic_inevitability}, partition operations at all these
sites had $\Delta G_{\mathrm{partition}} < 0$. They were not rare fluctuations;
they were thermodynamic necessities.
\end{corollary}

\begin{corollary}[Partition Rate]
\label{cor:partition_rate}
By transition state theory applied to partition operations with activation
barrier $\Delta G^{\ddagger}_{\mathrm{partition}}$:
\begin{equation}
  k_{\mathrm{partition}} = \frac{\kB T}{h} \exp\!\left(-\frac{\Delta G^{\ddagger}_{\mathrm{partition}}}{\kB T}\right).
  \label{eq:partition_rate}
\end{equation}
For mineral-surface partitions with $\Delta G^{\ddagger} \sim 20\,\mathrm{kJ\,mol}^{-1}$,
this gives $k_{\mathrm{partition}} \sim 10^7\,\mathrm{s}^{-1}$ at $T = 300\,\mathrm{K}$:
more than $10^7$ spontaneous partition operations per second per mineral surface site.
\end{corollary}

%------------------------------------------------------------------------------
\section{Theorem: Probability Separation}
\label{sec:probability_separation}
%------------------------------------------------------------------------------

\begin{theorem}[Probability Separation]
\label{thm:probability_separation}
The probability of a partition-first origin of life ($P_{\mathrm{partition}}$)
exceeds the probability of an information-first origin ($P_{\mathrm{information}}$)
by a factor greater than $10^{144}$:
\begin{equation}
  \frac{P_{\mathrm{partition}}}{P_{\mathrm{information}}} > 10^{144}.
  \label{eq:probability_ratio}
\end{equation}
This is a categorical distinction: information-first events are not merely
unlikely but are categorically impossible given the total number of trials
available in the observable universe.
\end{theorem}

\begin{proof}
We estimate each probability by counting the simultaneous independent conditions
that must be satisfied.

\medskip
\noindent\textbf{Information-first probability.}
A functional self-replicating RNA (the simplest information-first precursor)
requires, at minimum:

\begin{enumerate}[label=(\alph*)]
  \item \textbf{Specific sequence}: A ribozyme of 100 nucleotides capable of
        catalysing its own replication. From experimental surveys of sequence
        space, the fraction of random 100-mers with catalytic function is
        approximately $(1/20)^{100}$ for proteins (if protein-based catalysis
        is invoked) or $(1/4)^{100}$ for RNA, each corrected downward by the
        fraction of sequences that fold correctly ($\sim 10^{-10}$ of all
        sequences of sufficient length fold into compact functional structures).
        For the protein case (more commonly cited):
        \begin{equation}
          P_{\mathrm{seq}} \approx \left(\frac{1}{20}\right)^{100} \approx 10^{-130}.
          \label{eq:p_seq}
        \end{equation}
  \item \textbf{Catalytic function}: Even among compact-folding sequences,
        the fraction that are catalytically active for a specific reaction
        is empirically:
        \begin{equation}
          P_{\mathrm{cat}} \approx 10^{-10}.
          \label{eq:p_cat}
        \end{equation}
  \item \textbf{Membrane or compartment}: A concentration boundary sufficient
        to maintain the first replicator against dilution. Spontaneous vesicle
        formation requires amphiphiles with the right geometry (packing parameter
        $p \approx 1/3$), which form with probability:
        \begin{equation}
          P_{\mathrm{mem}} \approx 10^{-10}.
          \label{eq:p_mem}
        \end{equation}
\end{enumerate}

These conditions must be satisfied simultaneously in the same microscopic volume.
Assuming statistical independence (a conservative assumption that underestimates
the difficulty):
\begin{equation}
  P_{\mathrm{information}} \approx P_{\mathrm{seq}} \times P_{\mathrm{cat}} \times P_{\mathrm{mem}}
  \approx 10^{-130} \times 10^{-10} \times 10^{-10} = 10^{-150}.
  \label{eq:p_information}
\end{equation}

\medskip
\noindent\textbf{Partition-first probability.}
A partition-based precursor to life requires:

\begin{enumerate}[label=(\alph*)]
  \item \textbf{Chemical potential gradient}: Established wherever fluids of
        different chemical composition meet. The gradient $\Delta\mu_{\mathrm{H}^+}
        \approx 0.17\,\mathrm{eV}$ (3 pH unit proton gradient at alkaline
        hydrothermal vents) is the norm at Hadean vent fields:
        \begin{equation}
          P_{\mathrm{grad}} \approx 10^{-3}
          \label{eq:p_grad}
        \end{equation}
        (fraction of Hadean seafloor surface area with suitable vent conditions).
  \item \textbf{Semiconductor mineral surface}: Iron--sulfur minerals form
        spontaneously from Fe$^{2+}$ and H$_2$S in Hadean ocean water:
        \begin{equation}
          P_{\mathrm{surf}} \approx 10^{-2}
          \label{eq:p_surf}
        \end{equation}
        (fraction of vent fields with iron--sulfur mineralogy of sufficient area
        and semiconducting character).
  \item \textbf{Electron transport establishing partition}: The semiconducting
        surface in contact with the proton gradient spontaneously supports
        directed electron transport (driven by $\Delta G < 0$):
        \begin{equation}
          P_{\mathrm{et}} \approx 10^{-1}
          \label{eq:p_et}
        \end{equation}
        (fraction of suitable surfaces that develop organised redox chemistry).
\end{enumerate}

The joint probability is:
\begin{equation}
  P_{\mathrm{partition}} \approx 10^{-3} \times 10^{-2} \times 10^{-1} = 10^{-6}.
  \label{eq:p_partition}
\end{equation}

\medskip
\noindent\textbf{Ratio.}
\begin{equation}
  \frac{P_{\mathrm{partition}}}{P_{\mathrm{information}}}
  = \frac{10^{-6}}{10^{-150}} = 10^{144}. \qquad\blacksquare
  \label{eq:ratio_final}
\end{equation}
\end{proof}

\begin{remark}[Categorical Impossibility of Information-First]
The factor $10^{144}$ exceeds the total number of atoms in the observable
universe ($\sim 10^{80}$). Even granting every atom as a separate nucleotide,
every possible configuration tried at the quantum rate ($10^{13}\,\mathrm{s}^{-1}$)
for the full age of the universe ($4 \times 10^{17}\,\mathrm{s}$), the total
number of trials is $10^{80} \times 10^{13} \times 4 \times 10^{17} \approx 10^{111}$,
still $10^{39}$ times too few. The information-first scenario is not improbable;
it is physically impossible within our universe. The distinction between
partition-first and information-first is therefore categorical, not quantitative.
\end{remark}

\begin{table}[ht]
\centering
\caption{Comparison of origin-of-life scenarios by probability and key failure mode.}
\label{tab:origin_scenarios}
\begin{tabularx}{\textwidth}{l X X X}
\toprule
\textbf{Scenario} & \textbf{Requirements} & \textbf{Probability} & \textbf{Key failure} \\
\midrule
Partition-first
  & Redox gradient $+$ semiconductor mineral surface $+$ electron transport
  & $\sim 10^{-6}$ (probable)
  & None identified \\
\midrule
RNA world
  & Specific 100-mer sequence, catalytic fold, compartment
  & $\sim 10^{-150}$ (impossible)
  & Requires replication machinery it cannot pre-exist \\
\midrule
Metabolism-first
  & Autocatalytic coupled cycle ($\geq 5$ reactions), no inheritance
  & $\sim 10^{-30}$
  & No mechanism for reproducible specificity \\
\midrule
Lipid world
  & Amphiphile with packing parameter $p \approx 1/3$
  & $\sim 10^{-10}$
  & Compartmentalises without catalysing or encoding \\
\bottomrule
\end{tabularx}
\end{table}

%------------------------------------------------------------------------------
\section{Theorem: Partition Autocatalysis}
\label{sec:autocatalysis}
%------------------------------------------------------------------------------

Theorem~\ref{thm:thermodynamic_inevitability} establishes that the first
partition operation occurs. The following theorem establishes that it sustains
itself.

\begin{theorem}[Partition Autocatalysis]
\label{thm:partition_autocatalysis}
A partition operation that increases the local chemical potential gradient
at its own boundary is self-sustaining: the partitioning rate $\rho_{\mathrm{part}}$
is an increasing function of the partition depth $\Depth$,
\begin{equation}
  \frac{\partial \rho_{\mathrm{part}}}{\partial \Depth} > 0,
  \label{eq:autocatalysis_condition}
\end{equation}
creating a positive feedback loop.
\end{theorem}

\begin{proof}
Let $\Part$ be a partition with boundary $\partial\Part$ across which the
chemical potential is discontinuous by $\Delta\mu > 0$. The partitioning rate
follows from equation~\eqref{eq:dg_full}:
\begin{equation}
  \rho_{\mathrm{part}} = \rho_0 \exp\!\left(-\frac{\Delta G_{\mathrm{partition}}}{\kB T}\right)
  = \rho_0 \exp\!\left(\frac{\Delta\mu \cdot \Delta N}{\kB T}
    + \ln\!\frac{\Omega_{\mathrm{after}}}{\Omega_{\mathrm{before}}}\right).
  \label{eq:rho_part}
\end{equation}

The partition boundary $\partial\Part$ concentrates chemical species
near the boundary (by electrostatic focussing from the charge created at
$\partial\Part$, see Theorem~\ref{thm:charge_emergence}, Chapter~\ref{chap:charge_emergence}).
The local concentration $[c]_{\partial\Part}$ near the boundary satisfies:
\begin{equation}
  [c]_{\partial\Part} = [c]_{\mathrm{bulk}} \exp\!\left(\frac{q\phi_{\partial\Part}}{\kB T}\right),
  \label{eq:concentration_enhancement}
\end{equation}
where $\phi_{\partial\Part}$ is the electrostatic potential at the boundary
and $q$ is the charge of the species. The local chemical potential at the
boundary is:
\begin{equation}
  \mu_{\partial\Part} = \mu^0 + \kB T \ln[c]_{\partial\Part}
  = \mu^0 + \kB T \ln[c]_{\mathrm{bulk}} + q\phi_{\partial\Part}.
\end{equation}

The potential $\phi_{\partial\Part}$ grows with the charge accumulated at the
boundary. The charge accumulated is proportional to the partition depth $\Depth$
(each partition operation adds one charge layer to the boundary, by the Charge
Emergence Theorem). Therefore:
\begin{equation}
  \frac{\partial \phi_{\partial\Part}}{\partial \Depth} > 0,
  \quad\text{hence}\quad
  \frac{\partial \mu_{\partial\Part}}{\partial \Depth} > 0,
  \quad\text{hence}\quad
  \frac{\partial \Delta\mu}{\partial \Depth} > 0.
\end{equation}

Substituting into~\eqref{eq:rho_part}: since $\Delta\mu$ appears in the
exponential and $\partial\Delta\mu/\partial\Depth > 0$,
\begin{equation}
  \frac{\partial \rho_{\mathrm{part}}}{\partial \Depth}
  = \frac{\partial \rho_{\mathrm{part}}}{\partial \Delta\mu} \cdot \frac{\partial \Delta\mu}{\partial \Depth} > 0.
\end{equation}
The partitioning rate is an increasing function of partition depth: deeper
partitions drive faster further partitioning. This is autocatalysis. \qed
\end{proof}

\begin{corollary}[Self-Sustaining Partition Cascade: Proto-Metabolism]
\label{cor:partition_cascade}
By Theorem~\ref{thm:partition_autocatalysis}, a single spontaneous partition
operation on a suitable mineral surface initiates an autocatalytic cascade:
\begin{enumerate}[label=(\arabic*)]
  \item First partition drives directed electron transport (Theorem~\ref{thm:thermodynamic_inevitability}).
  \item Electron transport charges the partition boundary
        (Theorem~\ref{thm:charge_emergence}, Chapter~\ref{chap:charge_emergence}).
  \item Boundary charge concentrates reactants and intensifies $\Delta\mu$.
  \item Intensified $\Delta\mu$ drives further partition operations at rate
        $\rho_{\mathrm{part}} > \rho_0$.
  \item Network of partition operations reaches self-sustaining equilibrium.
\end{enumerate}
This self-sustaining partition network is proto-metabolism: a far-from-equilibrium
chemical network maintaining a flux of partition operations, without any enzyme.
\end{corollary}

%------------------------------------------------------------------------------
\section{Theorem: Chirality Selection}
\label{sec:chirality_selection}
%------------------------------------------------------------------------------

One of biochemistry's most striking regularities is universal homochirality:
all biological amino acids are L-configured, all biological sugars are
D-configured, and all double-helical DNA is right-handed. No exception is known
among the $\sim 500$ distinct chiral centres in the standard biochemical repertoire.
The origin of this universal selection has resisted explanation for decades:
standard chemistry produces racemic mixtures; an external chiral influence
(circularly polarised light, parity-violating weak force) provides at most
a tiny initial enantiomeric excess. The partition theory framework provides
a mechanistic account requiring none of these external sources.

\begin{theorem}[Chirality Selection via Electron Transport at Partition Boundaries]
\label{thm:chirality_selection}
Electron transport through chiral partition boundaries creates a chiral selection
bias for amino acid and sugar enantiomers by the following mechanism:
\begin{enumerate}[label=(\roman*)]
  \item Electrons carry spin angular momentum $\pm\hbar/2$.
  \item Transport through a chiral molecule transmits spin angular momentum
        (chirality-induced spin selectivity, CISS).
  \item The partition boundary geometry---set by the semiconductor mineral crystal
        surface---selects a preferred electron spin direction.
  \item Spin-selected electron transport preferentially stabilises L-amino acid
        and D-sugar configurations.
\end{enumerate}
The result is universal homochirality as a direct consequence of partition geometry.
No external chiral field is required.
\end{theorem}

\begin{proof}
\noindent\textbf{Step 1: CISS---spin-selective transmission through chiral molecules.}
Through a chiral molecule (a helix), electrons with spin parallel to the helix
axis are transmitted with efficiency $T_+$ and those with antiparallel spin with
efficiency $T_-$, with $T_+ \neq T_-$. This is the CISS effect, a direct
consequence of spin--orbit coupling in the helical potential. For a right-handed
helix, right-handed angular momentum (spin up) is preferentially transmitted
\citep{ray1999}:
\begin{equation}
  T_+(R\text{-helix}) > T_-(R\text{-helix}), \qquad
  T_+(L\text{-helix}) < T_-(L\text{-helix}).
  \label{eq:ciss}
\end{equation}

\noindent\textbf{Step 2: Surface spin--orbit coupling selects electron spin.}
The chiral mineral surface (iron--sulfur crystal termination) imposes a
Rashba-type spin--orbit interaction on electrons entering or leaving the surface:
\begin{equation}
  H_{\mathrm{Rashba}} = \lambda_{\mathrm{SO}}\,(\hat{n} \times \mathbf{k}) \cdot \boldsymbol{\sigma},
  \label{eq:rashba}
\end{equation}
where $\hat{n}$ is the surface normal, $\mathbf{k}$ is the electron wavevector,
$\boldsymbol{\sigma}$ is the spin operator, and $\lambda_{\mathrm{SO}}$ is the
Rashba coupling strength. For a right-handed crystal termination, equation~\eqref{eq:rashba}
prefers spin-up electrons in the direction of $\hat{n}$.

\noindent\textbf{Step 3: Spin selection biases enantiomer stabilisation.}
Electrons flowing across the chiral partition boundary (from reducing to
oxidising side, driven by $\Delta G < 0$) arrive at the surface with a preferred
spin determined by the surface crystal structure. By CISS (Step 1), these
spin-polarised electrons are preferentially transmitted through right-handed
molecular helices. L-amino acids ($\alpha$-helices are right-handed) and
D-sugars (in right-handed B-DNA) form right-handed structural motifs; therefore
they transmit the preferred spin more efficiently than their mirror images.

The energy difference between L- and D-amino acid adsorption at the chiral
partition boundary is:
\begin{equation}
  \Delta E_{\mathrm{chiral}} = 2g_s\mu_B B_{\mathrm{eff}},
  \label{eq:chiral_energy_gap}
\end{equation}
where $B_{\mathrm{eff}}$ is the effective magnetic field from the surface
spin--orbit coupling, and $g_s \approx 2$ is the electron $g$-factor.

\noindent\textbf{Step 4: Amplification to complete homochirality.}
At each polymerisation step, L-amino acid incorporation is preferred by the
Boltzmann factor $\exp(\Delta E_{\mathrm{chiral}}/\kB T) > 1$. Across $N$
incorporation events, the enantiomeric excess accumulates as:
\begin{equation}
  \mathrm{ee}(N) = \tanh\!\left(\frac{N \cdot \Delta E_{\mathrm{chiral}}}{\kB T}\right).
  \label{eq:ee_amplification}
\end{equation}
For a polymer of length $N = 1000$, even a small per-step preference
($\Delta E_{\mathrm{chiral}} / \kB T \approx 10^{-2}$) gives
$\mathrm{ee}(1000) = \tanh(10) \approx 1$: complete homochirality. The
partition geometry of the mineral surface selects chirality without invoking
any external chiral field. \qed
\end{proof}

\begin{remark}[On Universality]
Because all mineral surfaces of a given crystal type have the same chirality
(determined by the crystal symmetry group), all partition operations on that
mineral surface select the same enantiomer. This global coherence---all
partition-active sites biasing in the same direction---breaks the $\mathbb{Z}_2$
(L vs.\ D) symmetry globally. The biological homochirality of L-amino acids
and D-sugars reflects the chirality of the predominant Hadean iron--sulfur
mineral surface, not a cosmic accident.
\end{remark}

\begin{corollary}[D-Sugars and Right-Handed DNA]
\label{cor:dsugars_dna}
The same mechanism that selects L-amino acids selects D-ribose and
D-deoxyribose: the CISS transmission efficiency favours D-sugar configurations
at a right-handed chiral partition boundary. Right-handed B-form DNA follows
as the geometrically preferred duplex of two D-deoxyribose strands.
Universal homochirality across all three classes of biological chiral molecules
(amino acids, sugars, nucleic acid helices) follows from one mechanism applied
to one type of surface.
\end{corollary}

%------------------------------------------------------------------------------
\section{Prebiotic Chemistry at Low Temperature}
\label{sec:low_temp_prebiotic}
%------------------------------------------------------------------------------

Traditional abiogenesis scenarios favour warm-to-hot conditions: Darwin's ``warm
little pond,'' hydrothermal vents at $T > 100\,^\circ\mathrm{C}$, or volcanic
environments. This preference reflects the kinetic intuition that higher temperature
means faster reactions. Partition theory predicts a fundamentally different
controlling parameter.

\begin{proposition}[Partition-Temperature $T_\Part$]
\label{prop:partition_temperature}
Partition operations on semiconductor mineral surfaces are controlled by the
\emph{partition temperature} $T_\Part$, defined by the surface electrochemistry,
not by the ambient temperature $T_{\mathrm{ambient}}$:
\begin{equation}
  T_\Part = \frac{\Delta\mu \cdot \Delta N}{\kB \ln(\Omega_{\mathrm{after}} / \Omega_{\mathrm{before}})}.
  \label{eq:partition_temperature}
\end{equation}
For $T_{\mathrm{ambient}} < T_\Part$, partition operations are enthalpically
driven and remain spontaneous at arbitrarily low ambient temperature.
\end{proposition}

\begin{proof}
From equation~\eqref{eq:dg_full}, the partition operation is spontaneous if
$\Delta G_{\mathrm{partition}} < 0$, i.e., if:
\begin{equation}
  \Delta\mu \cdot \Delta N > -T_{\mathrm{ambient}}\,\kB \ln\!\frac{\Omega_{\mathrm{after}}}{\Omega_{\mathrm{before}}}.
  \label{eq:spontaneity_condition}
\end{equation}
Since $\ln(\Omega_{\mathrm{after}}/\Omega_{\mathrm{before}}) > 0$ and the right
side becomes more negative as $T_{\mathrm{ambient}}$ decreases, the condition
is \emph{easier} to satisfy at lower temperature. At $T_\Part$ (defined by
equality in equation~\eqref{eq:spontaneity_condition}), the partition operation
is at the margin of spontaneity. For all $T_{\mathrm{ambient}} < T_\Part$, the
partition operation is spontaneous. For $T_{\mathrm{ambient}} > T_\Part$, the
entropic drive augments the enthalpic drive and the operation is even more
spontaneous. The key parameter is $\Delta\mu$ (surface electrochemistry), not $T_{\mathrm{ambient}}$.
\qed
\end{proof}

\begin{corollary}[Abiogenesis at Arbitrarily Low Ambient Temperature]
\label{cor:low_temp_abiogenesis}
Wherever $\Delta\mu > 0$ on a suitable mineral surface, partition operations
proceed regardless of $T_{\mathrm{ambient}}$. This predicts prebiotic
chemistry---and potentially life---in cold molecular clouds and on the surfaces
of interstellar dust grains, consistent with:
\begin{itemize}
  \item Complex organic molecules (glycolaldehyde, methyl formate, acetamide,
        glycine precursors) detected in the Sagittarius B2 molecular cloud
        complex at $T \approx 10$--$80\,\mathrm{K}$ \citep{belloche2008,hollis2000}.
  \item Amino acid precursors and simple sugars in carbonaceous chondrites,
        formed in the cold outer solar system.
  \item Laboratory synthesis of complex organics by UV irradiation of cold
        ($T \approx 10\,\mathrm{K}$) ices of H$_2$O, CH$_3$OH, NH$_3$, and CO,
        where the driving force is photon energy ($\Delta\mu$ from UV), not
        thermal activation.
\end{itemize}
The relevant parameter for prebiotic chemistry is the surface chemical potential
gradient $\Delta\mu$, which can be maintained by UV photons, cosmic-ray
ionisation, or mineral redox chemistry at any ambient temperature.
\end{corollary}

%------------------------------------------------------------------------------
\section{Membranes as Partition Scaffolding}
\label{sec:membranes_partition}
%------------------------------------------------------------------------------

The traditional view holds that lipid membranes evolved to \emph{compartmentalise}
chemistry---to enclose a growing molecular system within a protective boundary.
In this account, membranes serve the genetic system and are selected for that
service. The partition theory account is different in both sequence and causation.

\begin{theorem}[Membranes as Physical Partition Structures]
\label{thm:membranes_partition}
Lipid bilayer membranes are physical instantiations of partition operations.
They form because they \emph{are} partition structures, not because they serve
the genetic system. Compartmentalisation is a geometric consequence of a closed
partition boundary in three dimensions, not the function of the membrane.
\end{theorem}

\begin{proof}
An amphiphilic molecule has a hydrophilic head group and a hydrophobic tail.
The hydrophobic effect imposes a chemical partition: contact between the
hydrophobic tail and water is thermodynamically penalised by approximately
$\Delta G_{\mathrm{exposure}} \approx 40\,\mathrm{J\,m}^{-2}$ of hydrophobic
surface area. This is a partition operation: the system \emph{partitions} into
hydrophilic and hydrophobic chemical state spaces with $\Delta G_{\mathrm{partition}} < 0$
(by Theorem~\ref{thm:thermodynamic_inevitability}, with $\Delta\mu$ the chemical
potential difference between aqueous and hydrocarbon environments).

The lipid bilayer is the unique geometric structure that minimises total
hydrophobic exposure in aqueous solution: tails pack against tails (forming
the interior), heads face water (forming both surfaces). In the partition
framework, the bilayer is the physical realisation of the chemical partition
$\Cat = \Cat_{\mathrm{hydrophobic}} \cup \Cat_{\mathrm{hydrophilic}}$, with
$\partial\Cat$ being the bilayer interfaces. The bilayer satisfies all conditions
of Definition~\ref{def:chemical_partition}:
\begin{itemize}
  \item \textit{Selectivity}: charged molecules cannot cross the hydrophobic
        interior without energy $\Delta G_{\mathrm{cross}} \approx 40$--$80\,\mathrm{kJ\,mol}^{-1}$.
  \item \textit{Stability}: maintained by the hydrophobic effect at all
        biologically relevant temperatures.
  \item \textit{Gradient coupling}: the membrane supports ion and proton
        gradients ($\Delta\mu_{\mathrm{ion}}$ across the bilayer is the
        proton-motive force driving ATP synthesis).
\end{itemize}

Compartmentalisation is a geometric corollary: closing a partition boundary
in three-dimensional space necessarily encloses a volume. The enclosed volume
is a consequence of the closed partition surface, not the reason for forming it.
The membrane did not evolve to compartmentalise; it formed because it is
thermodynamically favoured, and compartmentalisation followed geometrically.
\qed
\end{proof}

\begin{corollary}[Membrane Potential as Partition-Generated Charge]
\label{cor:membrane_potential}
The resting membrane potential ($\sim -70\,\mathrm{mV}$ in neurons,
$\sim -120\,\mathrm{mV}$ in mitochondria) is not the result of pre-existing
charges being mechanically separated by the membrane. By the Charge Emergence
Theorem (Chapter~\ref{chap:charge_emergence}), every partition boundary
carries charge proportional to the partition depth gradient across it. The lipid
bilayer, as a partition boundary with asymmetric partition depth across its two
leaflets (due to asymmetric lipid composition and integral protein distribution),
necessarily generates a potential difference. The membrane potential is
thermodynamically predicted by the partition framework; it does not require
pre-existing ions being pumped into position.
\end{corollary}

%------------------------------------------------------------------------------
\section{Resolution of Orgel's Paradox}
\label{sec:orgel_resolution}
%------------------------------------------------------------------------------

The four theorems of Sections~\ref{sec:thermodynamic_inevitability}--\ref{sec:membranes_partition}
converge on a clean resolution. We now state it.

\begin{theorem}[Partition-First Resolution of Orgel's Paradox]
\label{thm:orgel_resolution}
Partition operations on semiconductor mineral surfaces with chemical potential
gradients simultaneously and independently generate:
\begin{enumerate}[label=(\roman*)]
  \item Information: the partition structure $\Part$ stores $\Depth$ bits of
        categorical information (Theorem~\ref{thm:depth_entropy},
        Chapter~\ref{chap:bounded_phase_space}).
  \item Catalysis: partition topology change is the substrate of catalytic
        function (Chapter~\ref{chap:enzymes_topology}).
  \item Metabolism: the self-sustaining partition cascade maintains a
        far-from-equilibrium chemical network (Corollary~\ref{cor:partition_cascade}).
\end{enumerate}
None of the three requires the others to exist first. Orgel's paradox is
resolved because the apparent circularity arises from treating genes, enzymes,
and metabolism as irreducible primitives rather than recognising them as three
coordinate descriptions of the same underlying operation.
\end{theorem}

\begin{proof}
\noindent\textbf{(i) Information from partition structure.}
A partition $\Part$ of depth $\Depth$ encodes $\Depth$ bits of categorical
information (Theorem~\ref{thm:depth_entropy}, Chapter~\ref{chap:bounded_phase_space}:
$S_\Part = \kB \Depth \ln b$). On a mineral surface, the partition structure
records which surface sites are in which chemical states. This is information
storage without genetic molecules. The partition \emph{is} the proto-genome.

\noindent\textbf{(ii) Catalysis from partition topology change.}
A catalyst reduces the categorical path length from substrate to product
partition state: $\Delta\Depth_{\mathrm{path}}^{\mathrm{cat}} < \Delta\Depth_{\mathrm{path}}^{\mathrm{uncat}}$.
The mineral surface provides alternative partition topologies that reduce this
path length for specific substrate structures. Surface-driven partition topology
change is catalysis without any protein enzyme.

\noindent\textbf{(iii) Metabolism from autocatalytic partition cascade.}
By Theorem~\ref{thm:partition_autocatalysis} and Corollary~\ref{cor:partition_cascade},
partition operations on a gradient-bearing surface constitute a self-sustaining
chemical network: they consume gradient energy and maintain far-from-equilibrium
chemical composition. This is metabolism in its essential sense---energy
transduction maintaining non-equilibrium states---without any enzymatic machinery.

\noindent\textbf{Simultaneity.} All three arise from the identical physical
process: partition operations on a mineral surface with $\Delta\mu > 0$. They
therefore arise simultaneously, not sequentially. No circularity exists because
none of the three is the \emph{cause} of the partition operation; the partition
operation is the common cause of all three. \qed
\end{proof}

%------------------------------------------------------------------------------
\section{A Quantitative Timeline of Early Earth}
\label{sec:early_earth_timeline}
%------------------------------------------------------------------------------

\begin{enumerate}[label=\textbf{\arabic*.}]
  \item \textbf{4.4 Gya: Semiconductor mineral surfaces form.}
        Zircon ages establish continental crust by $4.4\,\mathrm{Gya}$.
        Hadean ocean chemistry (Fe$^{2+}$-rich, H$_2$S-bearing, pH $\approx 5$--$6$)
        deposits iron--sulfur minerals at hydrothermal sites. Partition operations
        begin spontaneously ($\Delta G_{\mathrm{partition}} < 0$ by
        Theorem~\ref{thm:thermodynamic_inevitability}).

  \item \textbf{4.4--4.2 Gya: Electron transport and charge separation.}
        Semiconductor surfaces couple to the proton gradient between alkaline
        vent fluids (pH $\approx 9$--$11$) and the acidic Hadean ocean.
        Directed electron transport reduces \ce{CO2} and \ce{N2} to organic
        molecules. Charge separation at the partition boundary generates the
        first electrochemical potential (Charge Emergence Theorem,
        Chapter~\ref{chap:charge_emergence}).

  \item \textbf{Concurrent: Chirality selection.}
        Electron transport through the right-handed chiral FeS crystal
        termination preferentially stabilises L-amino acid and D-sugar
        configurations (Theorem~\ref{thm:chirality_selection}). Homochirality
        is selected before the first genetic polymer exists.

  \item \textbf{4.2--4.0 Gya: Autocatalytic partition cascade (proto-metabolism).}
        By Theorem~\ref{thm:partition_autocatalysis}, first partition operations
        deepen and accelerate further partitioning. A self-sustaining network
        of partition operations forms---the first metabolism, requiring no
        enzyme, only mineral redox chemistry and the chemical gradient.

  \item \textbf{$\sim$4.0 Gya: Nucleotide stabilisation at partition boundaries.}
        Nucleotide ring systems (purines, pyrimidines) concentrate at partition
        boundaries through aromatic $\pi$-stacking interactions with the mineral
        surface. Nucleoside phosphates form as thermodynamically stabilised
        partition structures.

  \item \textbf{$\sim$4.0--3.8 Gya: Replication as partition copying.}
        Template-directed synthesis of complementary nucleotides is partition
        copying: each nucleotide at the boundary selects its Watson--Crick
        complement by the same partition completion criterion established in
        Chapter~\ref{chap:charge_emergence}. No polymerase enzyme is required
        for the first replication; the mineral surface and the geometry of
        the existing nucleotide chain are sufficient.

  \item \textbf{$\sim$3.8 Gya: Enzyme evolution as partition topology optimisation.}
        Polypeptides that provide better partition topology (reducing categorical
        path length more efficiently than the mineral surface) are selectively
        replicated. Enzymes evolve as optimised partition topologies, replacing
        the mineral surface with a more flexible and evolvable partition medium.
        Protein-based biology begins.
\end{enumerate}

\begin{remark}[Geological Consistency]
The oldest geological evidence for life---carbon isotope fractionation consistent
with biological metabolism in Isua supracrustal rocks---dates to approximately
$3.7$--$3.8\,\mathrm{Gya}$ \citep{ohtomo2014}. The timeline above places
a complete partition-based metabolism by $\sim 4.0\,\mathrm{Gya}$ and
enzyme-based biology by $\sim 3.8\,\mathrm{Gya}$. The $\sim 0.4\,\mathrm{Gyr}$
window between stable crust ($4.4\,\mathrm{Gya}$) and first geological evidence
($\sim 4.0\,\mathrm{Gya}$) is amply sufficient for steps 1--4 of the cascade,
without requiring any astronomically improbable event.
\end{remark}

%------------------------------------------------------------------------------
\section{Synthesis: Information, Catalysis, and Metabolism as Partition Triad}
\label{sec:partition_triad}
%------------------------------------------------------------------------------

The resolution of this chapter can be expressed compactly.

\begin{proposition}[The Partition Triad]
\label{prop:partition_triad}
Information, catalysis, and metabolism are three coordinate descriptions of
partition operations, each corresponding to a different quantity derived from
the partition:

\begin{center}
\begin{tabular}{lll}
\toprule
\textbf{Partition quantity} & \textbf{Coordinate description} & \textbf{Biological manifestation} \\
\midrule
Depth $\Depth$ & stored categorical distinctions & genetic information \\
Path length $\Delta\Depth_{\mathrm{path}}$ & change in partition topology & enzymatic catalysis \\
Rate $d\Depth/dt$ & maintenance of partition cascade & metabolic energy flux \\
\bottomrule
\end{tabular}
\end{center}
All three are derivable from the partition axioms of
Chapter~\ref{chap:bounded_phase_space} and arise together wherever partition
operations occur.
\end{proposition}

\begin{proof}
\noindent\textbf{Information as depth.} By Theorem~\ref{thm:depth_entropy}
of Chapter~\ref{chap:bounded_phase_space}, $S = \kB \Depth \ln b$, and by
the Shannon--Boltzmann identification, $\Depth$ is the information content in
bits. A genetic sequence is a partition of nucleotide sequence space; its
information content is its partition depth.

\noindent\textbf{Catalysis as path length reduction.} An enzyme reduces the
categorical path length from substrate partition state to product partition
state: $\Delta\Depth_{\mathrm{path}}^{\mathrm{cat}} < \Delta\Depth_{\mathrm{path}}^{\mathrm{uncat}}$
(Chapter~\ref{chap:enzymes_topology}). A shorter categorical path corresponds
to a lower activation free energy. Catalysis is partition topology navigation.

\noindent\textbf{Metabolism as depth rate.} The autocatalytic partition cascade
(Corollary~\ref{cor:partition_cascade}) maintains $d\Depth/dt > 0$ for the
ensemble of partition operations, consuming gradient energy and producing
partition depth. This is metabolism in categorical coordinates: a sustained
positive rate of categorical distinction generation.

The three descriptions are not causally independent; they are projections of
the single partition process onto three different observable axes. \qed
\end{proof}

The Partition Triad (Proposition~\ref{prop:partition_triad}) is the formal
statement of this chapter's central claim: life is not an improbable accident
but a geometrically necessary consequence of partition operations in a chemical
gradient environment. The three components that appear circular in biochemical
terms are three coordinate descriptions of one underlying physical operation.

%------------------------------------------------------------------------------
\section{Connection to the Paradox Resolution Programme}
\label{sec:paradox_programme}
%------------------------------------------------------------------------------

The resolution pattern of this chapter mirrors that of Chapter~\ref{chap:maxwells_demon}
in a mathematically precise sense. In both cases, an apparent paradox arises
from apparent circular dependence between two (or more) quantities; in both
cases, the partition theory framework reveals that the quantities are not
independent primitives but projections of the same more fundamental operation.

In Maxwell's Demon: the paradox involves entropy (categorical, partition depth)
and kinetic energy (dynamical, velocity). The demon appears to use one to control
the other, but they are categorically orthogonal: $\partial S / \partial v = 0$
(Theorem~\ref{thm:categorical_kinetic_orthogonality}, Chapter~\ref{chap:maxwells_demon}).

In Orgel's paradox: the paradox involves information, catalysis, and metabolism.
Genes, enzymes, and metabolism each appear to require the others, but all three
are projections of the partition: they are not independent quantities but
coordinates of the partition framework (Proposition~\ref{prop:partition_triad}).

This pattern---dissolving apparent paradoxes by identifying the common categorical
substrate of apparently independent quantities---is the central method of
Part~\ref{part:paradoxes}. Chapters~\ref{chap:what_is_life} and~\ref{chap:petos_paradox}
will apply the same method to Schr\"{o}dinger's question of what distinguishes
living from non-living matter, and to Peto's paradox of why cancer incidence
does not scale with body size.

%=============================================================================
\section*{Chapter Summary}
\label{sec:ch6_summary}
%=============================================================================

\begin{chapterbox}[title=Chapter 6 --- Key Results Established]
\begin{enumerate}[label=\textbf{\arabic*.}]
  \item \textbf{Orgel's Circle} (Definition~\ref{def:orgels_circle}): The
        mutual dependencies among genes, enzymes, and metabolism constitute a
        genuine logical impossibility within any ``X-first'' molecular framework.
        RNA world, metabolism-first, and lipid world each resolve one arc of
        the circle while restoring it elsewhere.

  \item \textbf{Chemical Partition Operations} (Definition~\ref{def:chemical_partition}):
        Selective, stable, gradient-coupled division of chemical state space.
        Iron--sulfur mineral surfaces at Hadean hydrothermal vents are natural
        partition structures (Example~\ref{ex:iron_sulfur_partition}).

  \item \textbf{Thermodynamic Inevitability} (Theorem~\ref{thm:thermodynamic_inevitability}):
        $\Delta G_{\mathrm{partition}} < 0$ wherever $\Delta\mu > 0$. Partition
        operations are thermodynamically spontaneous in chemical gradient
        environments. On early Earth, they were not improbable accidents but
        unavoidable consequences of mineralogy.

  \item \textbf{Probability Separation} (Theorem~\ref{thm:probability_separation}):
        $P_{\mathrm{partition}} / P_{\mathrm{information}} = 10^{144}$.
        With $P_{\mathrm{partition}} \approx 10^{-6}$ and
        $P_{\mathrm{information}} \approx 10^{-150}$, the information-first
        scenario is categorically impossible given the total number of trials
        available in the observable universe.

  \item \textbf{Partition Autocatalysis} (Theorem~\ref{thm:partition_autocatalysis}):
        $\partial \rho_{\mathrm{part}} / \partial \Depth > 0$. Partition
        operations that intensify the gradient at their own boundary are
        self-sustaining. The resulting cascade (Corollary~\ref{cor:partition_cascade})
        is abiotic proto-metabolism---without enzymes.

  \item \textbf{Chirality Selection} (Theorem~\ref{thm:chirality_selection}):
        Electron transport through chiral partition boundaries selects
        L-amino acids and D-sugars via the CISS mechanism and Rashba
        spin--orbit coupling. Universal homochirality is a consequence of
        partition boundary geometry, not a cosmic accident. No external
        chiral field is required.

  \item \textbf{Low-Temperature Abiogenesis} (Proposition~\ref{prop:partition_temperature}
        and Corollary~\ref{cor:low_temp_abiogenesis}): Partition operations
        are controlled by $T_\Part$ (surface electrochemistry), not
        $T_{\mathrm{ambient}}$. Prebiotic chemistry is predicted at arbitrarily
        low ambient temperature, consistent with complex organics in cold
        molecular clouds at $T \approx 10$--$80\,\mathrm{K}$.

  \item \textbf{Membranes as Partition Structures} (Theorem~\ref{thm:membranes_partition}):
        Lipid bilayers are physical partition operations. They form because
        the hydrophobic partition has $\Delta G < 0$, not because they serve
        the genetic system. Compartmentalisation is geometric; membrane
        potential is partition-generated charge
        (Corollary~\ref{cor:membrane_potential}).

  \item \textbf{Orgel's Paradox Resolved} (Theorem~\ref{thm:orgel_resolution}):
        Information ($\Depth$), catalysis ($\Delta\Depth_{\mathrm{path}}$),
        and metabolism ($d\Depth/dt$) arise simultaneously and independently
        from partition operations on a mineral surface with $\Delta\mu > 0$.
        None requires the others. Orgel's circle dissolves.

  \item \textbf{Partition Triad} (Proposition~\ref{prop:partition_triad}):
        Information, catalysis, and metabolism are three coordinate descriptions
        of the partition. Life is the self-sustaining dynamics of a partition
        system in a gradient environment. This is not a metaphor; it is a
        coordinate transformation from biochemical language to the partition
        framework.
\end{enumerate}
\end{chapterbox}
